% Introduction -- drafted by loop0015 item 09 (2026-11-05), under the owner's
% plain-English rule of 2026-11-05 05:30: no formula beyond 2^n in this
% section, every term of art glossed in one clause on first use, every
% number with the section that regenerates it. The title's four clauses are
% made precise in the first paragraph; the symbols (the exponent, the
% constants) are introduced in Section 2 and used from there on. A blind
% reader's pass (one subagent, 2026-11-05) moved the glossary ahead of the
% summary and shortened the long sentences.
\section{Introduction}
\label{sec:intro}

In late October 2026, OpenAI's collection of machine-checked mathematics
added a result dated 25 October 2026 that it calls family 378,
\emph{The Exponential Time Hypothesis is False}
\citep{OAI:The-Exponential-Time-Hypothesis-is-False-October-25-2026}.
The Exponential Time Hypothesis, ETH, is a conjecture from 2001
\citep{ImpagliazzoPaturi2001}. It says that deciding whether a formula
with three variables per clause can be satisfied needs time that grows
exponentially in the number of variables. The result comes with a
67-page manuscript and a proof checked by the Lean proof assistant, a
program that verifies every step of a proof and rejects a proof with a
gap. The verified theorem does refute the hypothesis. The title of this
paper says what else is true of it. Its sentence has four clauses.
First, what is verified is a \emph{guarantee}, not a running time. The
theorem says that a machine halts within a stated bound; it says
nothing about how long the machine actually takes on any input. Second,
the bound contains constants that the theorem only says \emph{exist}.
No value of any of them is computed anywhere. Where the proof does
write a value down, the constant in the exponent is at least 257
(Section~\ref{sec:statement}). Third, the \emph{exponent}, the quantity
that two is raised to in the bound, is a fixed function of the number
of variables. With every constant set to one, it is larger than the
number of variables at every size from 50 variables up to a number with
2,391 digits (Section~\ref{sec:statement}, Figure~\ref{fig:central}).
Over that whole range the verified bound is above $2^n$, the cost of
trying every one of the $2^n$ assignments of $n$ variables. So the
verified guarantee is weaker, at every practical size, than the
guarantee of the method everyone already has. Fourth, the manuscript
states two sharper exponents, which would make a different sentence.
Neither is in the Lean library, and the manuscript itself calls the
verified one ``the weaker exponent'' (Section~\ref{sec:beyond}).

This paper is a careful reading by people who checked, not a debunking.
We read the verified statement line by line. We inspected the proof's
source. We rebuilt the proof from source in isolation and re-checked it
with a second proof checker that shares no code with Lean's own. We
found no flaw: the statement means what the manuscript's verified form
says, and the proof rests on the three standard axioms, the assumptions
every ordinary Lean proof makes, and nothing else. Every sentence that
could read as ``wrong'' or ``overclaimed'' is replaced by what the Lean
statement and the manuscript actually say, with the file and line, or
the page and line, beside it. Nothing is quoted from memory. Every
number has a script that regenerates it, named where the number
appears.

\subsection{The words we use}
\label{sec:intro-words}

Each section glosses its own terms where they first appear, so that a
reader may enter anywhere. The ones the whole paper needs are these,
each in a sentence or two.
A \emph{formula} is a list of clauses, each clause an ``or'' of a few
variables or negated variables. A choice of true or false for every
variable \emph{satisfies} the formula when it makes every clause true.
A \emph{$k$-CNF} is a formula with at most $k$ variables per clause,
and $k$ is called the \emph{clause width}; 3-SAT is the problem of
deciding whether a 3-CNF can be satisfied.
\emph{Trying all $2^n$ assignments} is the method of checking every one
of the $2^n$ ways to set $n$ variables; its cost is the yardstick in the
title.
\emph{ETH} is the conjecture that 3-SAT needs time exponential in $n$.
\emph{SETH}, its strong form, says that as $k$ grows no algorithm for
$k$-SAT does much better than trying all assignments. Several variants
exist, and Section~\ref{sec:consequences} names each.
A \emph{guarantee} is a bound a theorem promises; a \emph{running time}
is what a clock measures; the two are compared only with their own
kind.
\emph{Lean} is the proof assistant. Its \emph{kernel} is the small core
that does the final checking, and ``verified'' in this paper always
means ``accepted by the kernel''. \emph{Mathlib} is Lean's community
library of mathematics.
The \emph{comparator} is OpenAI's checking tool. It takes a
\emph{challenge file}, which states a theorem with the proof left
blank, and a \emph{solution}, which proves it. It checks that both
state the same theorem and that the proof uses only the three standard
axioms.
\emph{nanoda} is a second, independent checker for Lean proofs, written
in another language. A \emph{replay} is a rebuild of a proof from
source on our own machine, inside a \emph{sandbox}, an isolated
environment that sees nothing else on the machine, at the \emph{pinned}
Lean version, the exact version the repository names.
A \emph{SAT solver} is the program a practitioner runs to decide a
formula, and a \emph{model counter} the one that counts its satisfying
assignments.

\subsection{What this paper does}
\label{sec:intro-summary}

In one paragraph. The verified statement (Section~\ref{sec:statement})
says that for each fixed clause width there exist a finite machine of
the textbook kind, a Turing machine with several tapes, and five
positive constants. On every formula of that width the machine halts
with the exact number of satisfying assignments written on a tape. It
does so within a bound: a constant, times a power of the input length,
times two raised to a constant times the exponent. The exponent is
written out in Section~\ref{sec:statement}. Its ratio to the number of
variables rises above one at 50, peaks at about 2.17 near a number with
24 digits, and falls back below one only past a number with 2,391
digits. The manuscript claims more than the statement
(Section~\ref{sec:beyond}): two sharper exponents; a bound on the size
of logic circuits that depends on one of them; and a chapter of proofs
by a technique of its own, tensor contraction. We tag every claim
verified, supported on paper by the verified statement, or paper-only.
We inspected the proof (Section~\ref{sec:proof}): all 3,091 source
files it depends on. They contain no axiom, no omitted proof, no step
that trusts a compiled program, and no patch to Mathlib, and the
statement they prove is the one the comparator posed. We replayed it
(Section~\ref{sec:replay}). The 3,091 files built from source in a
sandbox at the pinned Lean version in 33 minutes; the kernel reported
the three standard axioms; and nanoda re-checked 95,176 declarations,
the theorem and everything it depends on, with no error. We say what
follows (Section~\ref{sec:consequences}). The statement alone refutes
ETH, SETH, the gap form, the form for guessing (nondeterministic)
machines, the counting form, and the randomised forms of ETH, SETH and
the gap form. A lower bound that was proved \emph{from} one of them is
an ``if \ldots{} then'' whose ``if'' is now false: it derives nothing,
and the statement it concluded is not itself refuted. We say what it
changes for a practitioner (Section~\ref{sec:practitioner}): nothing on
any machine. We show that by timing an ordinary SAT solver and two model
counters on random formulas of up to 500 variables and drawing the two
guarantees beside the measurements. We restate the verified sentence in
our own Lean library under one name, with its axioms audited
(Section~\ref{sec:lean}). And we place the paper beside what the
collection, the hypotheses' authors, a published human audit, and the
two checking tools say about themselves (Section~\ref{sec:related}).

\subsection{What is claimed, and what is not}
\label{sec:intro-claims}

Claimed, each with the section that establishes it:
\begin{enumerate}
\item The verified statement is a running-time \emph{guarantee} with
  constants that exist and are never computed, for each fixed clause
  width, with no lower limit on the number of variables
  (Section~\ref{sec:statement}).
\item With every constant set to one, the verified exponent exceeds the
  number of variables at every size from 50 up to a number with 2,391
  digits, so the verified bound is above $2^n$ there. The true constants
  are unknown and at least 257 in the proof on record, so with them the
  bound is higher still (Section~\ref{sec:statement}).
\item The manuscript's two sharper exponents, and the circuit bound that
  rests on one of them, are not in the Lean library
  (Section~\ref{sec:beyond}).
\item The proof's source shows none of the known ways of weakening a
  verified proof, and its statement is the challenge's
  (Section~\ref{sec:proof}). It builds from source at the pinned
  version with the three standard axioms, and a second checker agrees
  (Section~\ref{sec:replay}).
\item The verified statement refutes ETH, SETH, their gap, guessing and
  counting forms, and the randomised forms of ETH, SETH and the gap
  form, by the parts of the statement named in each case. A lower bound
  derived from any of them loses its premise and nothing more
  (Section~\ref{sec:consequences}).
\item On random formulas, the solvers a practitioner has are many
  orders of magnitude below $2^n$ at every size measured, and the
  verified guarantee is above $2^n$ from 50 variables on. The verified
  result changes nothing that runs (Section~\ref{sec:practitioner}).
\end{enumerate}
Not claimed:
\begin{enumerate}
\item No running time of the verified machine, in either direction. No
  implementation of it exists; the title compares two guarantees.
\item No flaw in the statement, the proof, the Lean kernel, Mathlib, or
  the collection's checking tool. We found none and assert none.
\item No value of any constant beyond what the proof writes down, and
  nothing about what a different proof of the same statement could
  achieve.
\item No judgement on the manuscript's paper-only proofs, which were
  located and quoted, not read for correctness.
\item No refutation of any lower bound's conclusion, for any problem,
  and nothing about any specific problem's complexity beyond what the
  manuscript's own remarks state.
\item Nothing about any problem outside satisfiability. This paper is
  about the limitation of the ETH refutation and nothing else.
\end{enumerate}

\subsection{How the paper is organised}
\label{sec:intro-organisation}

Section~\ref{sec:statement} reads the statement.
Section~\ref{sec:beyond} draws the line between it and the manuscript.
Sections~\ref{sec:proof} and \ref{sec:replay} inspect and replay the
proof. Section~\ref{sec:consequences} says what follows.
Section~\ref{sec:practitioner} is for the practitioner.
Section~\ref{sec:lean} is our own Lean. Section~\ref{sec:related} is
related work. Section~\ref{sec:conclusion} concludes and lists the
questions this reading left for its author to decide. Every section
ends with a list headed ``what this does not claim''.
Appendix~\ref{app:lean} lists every Lean name the paper cites.
Appendix~\ref{app:tech} holds the technical detail behind the plain
text of the body: the definitions with their line numbers, the theorem
with its quantifiers, the spine of the proof. The OpenAI repository is
read throughout at one commit, \texttt{4eaee875}, as untrusted data:
with text tools only, never built or run in place. The one build is the
sandboxed replay.

\subsection{What this does not claim}
\label{sec:intro-not}

\begin{itemize}
\item \textbf{Nothing beyond the sections it summarises.} Every
  sentence of this introduction is a sentence of a later section said
  shorter; where the two differ, the later section, with its file and
  line, is the one that counts.
\item \textbf{Nothing about the result's importance.} Whether the
  refutation matters is not a question this paper asks; what the
  verified statement says, and at what sizes, is.
\item \textbf{No prior-art claim.} ``We found no earlier statement'' of
  the title's observation is bounded to the six texts
  Section~\ref{sec:related} names; the literature at large was not
  searched.
\end{itemize}
